\documentclass[11pt,a4paper]{article}
\usepackage{ctex}
\usepackage{amsmath,amssymb}
\usepackage{hyperref}

\title{\textbf{良知驱动的全息元认知：五维心智模型}\\
\large —— 一种可证伪的科学假设}
\author{\textbf{孟凡淳 (Grit Meng)}\\
前联想集团全球供应链集成计划系统 (IPS) 首席架构师\\
\texttt{gritmeng@outlook.com}}
\date{2026年9月}

\begin{document}
\maketitle

\begin{abstract}
本专著继承钱学森先生“开放复杂巨系统（OCGS）”与“大成智慧（Metasynthesis）”理论框架，以曹龙兵教授的非独立同分布（Non-IID）理论为唯一本体论源头，首次提出了针对 $\mathcal{O}(N!)$ 复杂巨系统的微观碳基 L0 “五维心智认知操作系统”。本书通过非平衡态热力学、微分拓扑与控制论，严密建立了具身良知 ($E_{\text{supp}}$)、残差敏感 ($\nabla R$)、感性缓存 ($\mathcal{S}_{\text{buff}}$)、流体智力 ($\hat{\mathcal{A}}_{\text{fluid}}$) 与全息元认知 ($\hat{\Omega}_{\text{meta}}$) 的算子闭环，完成了与大脑四大神经网络（DMN、SN、FPN、rPFC/BA10）的微分同构映射，提出了硅基 LLM 硬件熔断与硅基痛感模拟路径，并声明了三大严苛的可证伪性科学准则。

\vspace{0.5em}
\noindent\textbf{关键词}：五维心智模型；全息元认知；具身良知；紧支撑势垒；非独立同分布 (Non-IID)；开放复杂巨系统；硬件级熔断；可证伪性准则；碳硅协同
\end{abstract}
\end{document}
